\scanpage{82}
\mspage{17}
\[
(59)\qquad \rho_{n-1}\rho_{n-2}\cdots\rho_1\rho_0=1,
\]
et on trouve en fait
\[
\textcircled{60}\qquad
\Pi_P\simeq\Pi_{P_0}\simeq
\{\rho_0,\rho_1,\ldots,\rho_{n-1}\mid
\rho_{n-1}\rho_{n-2}\cdots\rho_1\rho_0=1\},
\]
où $P_0$ désigne le polygône combinatoire-type d'ordre $n$, i.e. muni d'un épinglage par un repère.

Venons-en à $\Pi^\tau$, on introduit les
\[
(61)\qquad \tau_a=(\lambda_a,\tau)\in\Pi_P^\tau\qquad(a\in A)
\]
on trouve que
\[
\textcircled{62}\qquad \Pi^\tau\simeq\{(\tau_a)_{a\in A}\mid\tau_a^2=1\},
\]
les relations entre les $\tau_a$ et les $\rho_s$ étant données par
\[
(63)\qquad \rho_s=\tau_a\tau_{a'},\qquad
\begin{gathered}
\text{où }a,a'\text{ sont les deux arêtes incidentes}\\
\text{à }s,\ \deleted{\unsure{tel}}\text{ avec }\deleted{\unsure{que}}a=a(s,\omega)\text{ i.e.}\\
a=\omega(a')\quad\text{i.e. }a'=\omega'(a).
\end{gathered}
\]
Si on choisit \deleted{\unsure{aussi}} une arête origine $a_0$
\[
a_0=a(s_0,\omega),\quad\text{donc }(s_0,a_0)\text{ repère orienté}
\]
en écrivant
\[
(64)\qquad \tau_i=\tau_{\omega^i a_0}\qquad(i\in\mathbb Z/n\mathbb Z)
\]
on trouve donc, avec la notation (58)
\[
(65)\qquad \rho_i=\tau_i\tau_{i-1}
\]
(de sorte que la relation (59) est conséquence des relations $\tau_i^2=1$). L'écriture (62) devient
\[
(62\ \text{bis})\qquad
\Pi^\tau\simeq\{\tau_0,\ldots,\tau_{n-1}\mid
\tau_0^2=\tau_1^2\cdots=\tau_{n-1}^2=1\}.
\]
\margin{Attention, le $\tau_{a_0}$ noté ici $\tau_0$, est noté $\tau_\infty$ dans \unsure{les textes antérieurs}, pour $n=3$ (car $a_0$ est ``opposé'' au sommet $R_\infty$...).}

Considérons maintenant le stabilisateur de $P_\omega$ dans $D_P\times\{1,\tau\}$, il est formé des
\[
(66)\qquad
\widetilde g^0\unsure{:}= (g,\tau^{\alpha(g)})
\qquad
\begin{gathered}
\alpha(g)\in\mathbb Z/2\mathbb Z\\
\alpha(g)=
\begin{cases}
0&\text{si }\operatorname{sg}(g)=1\\
1&\underline{\phantom{\text{si }\operatorname{sg}(g)}}=-1
\end{cases}
\end{gathered}
\]
\margin{81}
